% ═══════════════════════════════════════════════════════════════
% ARBEITSBLATT: parecer y conocer (DS 1)
% Template: Version E "Refined Hybrid"
% ═══════════════════════════════════════════════════════════════

\documentclass[10pt, a4paper]{article}

\usepackage[utf8]{inputenc}
\usepackage[T1]{fontenc}
\usepackage[ngerman]{babel}

% --- SCHRIFTEN ---
\usepackage{palatino}
\usepackage{helvet}
\newcommand{\sanstext}[1]{{\fontfamily{phv}\selectfont #1}}

\usepackage{microtype}
\usepackage[margin=1.4cm, top=1.6cm, bottom=1.3cm]{geometry}
\usepackage{enumitem}
\usepackage{tabularx}
\usepackage{booktabs}
\usepackage{array}
\usepackage{multicol}
\usepackage{tikz}
\usetikzlibrary{calc}

\usepackage{fancyhdr}
\usepackage{lastpage}
\usepackage{needspace}

\usepackage{xcolor}
\usepackage{tcolorbox}
\tcbuselibrary{skins}

% ═══════════════════════════════════════════════════════════════
% FARBEN
% ═══════════════════════════════════════════════════════════════

\definecolor{kopfzeile}{HTML}{1A1A1A}
\definecolor{akzent}{HTML}{C62828}
\definecolor{hellgrau}{HTML}{F2F2F2}
\definecolor{mittelgrau}{HTML}{777777}
\definecolor{dunkelgrau}{HTML}{444444}

% ═══════════════════════════════════════════════════════════════
% VARIABLEN
% ═══════════════════════════════════════════════════════════════

\newcommand{\thema}{parecer y conocer}
\newcommand{\klasse}{9}
\newcommand{\maxpunkte}{26}

% ═══════════════════════════════════════════════════════════════
% KOPF- UND FUSSZEILE
% ═══════════════════════════════════════════════════════════════

\pagestyle{fancy}
\fancyhf{}
\renewcommand{\headrulewidth}{0pt}
\renewcommand{\footrulewidth}{0pt}
\cfoot{\sanstext{\footnotesize\textcolor{mittelgrau}{\thepage\,/\,\pageref{LastPage}}}}

% ═══════════════════════════════════════════════════════════════
% BOXEN
% ═══════════════════════════════════════════════════════════════

\newtcolorbox{grammatikbox}{
    colback=hellgrau,
    colframe=hellgrau,
    boxrule=0pt,
    arc=0pt,
    left=8pt, right=8pt, top=8pt, bottom=6pt,
    borderline north={1.5pt}{0pt}{dunkelgrau}
}

\newtcolorbox{vokabelbox}{
    colback=white,
    colframe=mittelgrau,
    boxrule=0.5pt,
    arc=0pt,
    left=6pt, right=6pt, top=5pt, bottom=5pt
}

\newcommand{\regel}[1]{%
    \tikz[baseline=(X.base)]\node[fill=akzent, text=white,
    font=\sanstext\scriptsize\bfseries, inner sep=2pt, rounded corners=1pt](X){#1};%
}

% ═══════════════════════════════════════════════════════════════
% BEFEHLE
% ═══════════════════════════════════════════════════════════════

\newcommand{\es}[1]{\textbf{#1}}
\newcommand{\de}[1]{{\small\textit{\textcolor{mittelgrau}{#1}}}}
\newcommand{\stamm}[1]{\textcolor{akzent}{#1}}
\newcommand{\luecke}[1]{\rule{#1}{0.4pt}}
\newcommand{\punktebox}[1]{\hfill\sanstext{\small[\_\_/#1]}}

\newcommand{\aufgabe}[2]{%
    \needspace{4\baselineskip}%
    \vspace{12pt}%
    \noindent\sanstext{\textbf{#1}\quad#2}%
    \vspace{6pt}%
    \nopagebreak[4]%
}

\newlist{teil}{enumerate}{1}
\setlist[teil]{label=\alph*), leftmargin=1.8em, itemsep=3pt, parsep=0pt, topsep=3pt}

\setlength{\parindent}{0pt}
\setlength{\parskip}{4pt}
\setlength{\columnsep}{24pt}

% ═══════════════════════════════════════════════════════════════
% DOKUMENT
% ═══════════════════════════════════════════════════════════════

\begin{document}

% ═══════════════════════════════════════════════════════════════
% HEADER
% ═══════════════════════════════════════════════════════════════

\noindent
\begin{tikzpicture}[remember picture, overlay]
    \fill[kopfzeile] (current page.north west) rectangle +(\paperwidth, -1cm);
    \node[anchor=west, text=white, font=\sanstext\large\bfseries]
        at ([xshift=1.4cm, yshift=-0.5cm]current page.north west) {\thema};
    \node[anchor=east, text=white, font=\sanstext\small]
        at ([xshift=-1.4cm, yshift=-0.5cm]current page.north east)
        {Spanisch\,·\,Klasse \klasse\,·\,Arbeitsblatt};
\end{tikzpicture}

\vspace{0.5cm}

% --- Name/Datum ---
\noindent
\sanstext{\small\textbf{Name:}} \luecke{5cm} \hfill
\sanstext{\small\textbf{Datum:}} \luecke{2.5cm}

\vspace{6pt}

% ═══════════════════════════════════════════════════════════════
% EINLEITUNG
% ═══════════════════════════════════════════════════════════════

{\small
\textbf{¿Qué te parece?} -- Mit \es{parecer} kannst du deine Meinung sagen.
Mit \es{conocer} fragst du, ob jemand eine Person oder Sache kennt.
Beides brauchst du ständig im Alltag!}

\vspace{8pt}

% ═══════════════════════════════════════════════════════════════
% GRAMMATIK-ÜBERSICHT
% ═══════════════════════════════════════════════════════════════

\begin{grammatikbox}
\begin{multicols}{2}

% --- parecer ---
{\large\bfseries parecer} \de{scheinen, gefallen}

\vspace{4pt}

Konstruktion wie \es{gustar}:

\vspace{3pt}

{\small
\begin{tabular}{@{}ll@{}}
me parece(n) & \de{mir scheint} \\
te parece(n) & \de{dir scheint} \\
le parece(n) & \de{ihm/ihr scheint} \\
\end{tabular}}

\columnbreak

% --- conocer ---
{\large\bfseries conocer} \de{kennen}

\vspace{4pt}

\regel{yo → conozco} \quad {\small(c → zc)}

\vspace{6pt}

{\small
\begin{tabular}{@{}ll@{}}
\es{conocer a} & + Person \\
\es{conocer} & + Sache \\
\end{tabular}}

\end{multicols}

\vspace{4pt}

\begin{center}
{\small
\regel{e→\stamm{ie}} \es{qu\stamm{ie}ro, qu\stamm{ie}res, qu\stamm{ie}re}, queremos, queréis, \es{qu\stamm{ie}ren}
\hspace{6pt}
\regel{o→\stamm{ue}} \es{p\stamm{ue}do, p\stamm{ue}des, p\stamm{ue}de}, podemos, podéis, \es{p\stamm{ue}den}}
\end{center}
\end{grammatikbox}

\vspace{6pt}

% ═══════════════════════════════════════════════════════════════
% AUFGABEN (zweispaltig)
% ═══════════════════════════════════════════════════════════════

\begin{multicols}{2}

% --- AUFGABE 1 ---
\aufgabe{1}{Ergänze \es{parecer}.} \punktebox{6}

\begin{teil}
    \item \es{¿Qué te} \luecke{2cm} \es{los deberes?}\\
          --- \es{Me} \luecke{2cm} \es{difíciles.}
    \item \es{¿Qué os} \luecke{2cm} \es{la película?}\\
          --- \es{Nos} \luecke{2cm} \es{muy divertida.}
    \item \es{A Clara le} \luecke{2cm} \es{que Gonzalo es un poco raro.}
    \item \es{¿Qué te} \luecke{2cm} \es{quedar a las ocho?}
\end{teil}

\vspace{6pt}

% --- AUFGABE 2 ---
\aufgabe{2}{Ergänze \es{conocer}.} \punktebox{7}

\begin{teil}
    \item \es{Yo no} \luecke{2cm} \es{a tu hermano.}
    \item \es{¿Tú} \luecke{2cm} \es{la calle Trinidad?}
    \item \es{Nosotros} \luecke{2cm} \es{a toda tu familia.}
    \item \es{Mis hermanos} \luecke{2cm} \es{a Alfonso.}
    \item \es{Clara todavía no} \luecke{2cm} \es{muy bien a Nicolás.}
    \item \es{¿Vosotros} \luecke{2cm} \es{a los amigos de Gonzalo?}
    \item \es{Yo} \luecke{2cm} \es{todas las calles del centro.}
\end{teil}

\columnbreak

% --- AUFGABE 3 ---
\aufgabe{3}{\es{a} oder keine Präposition?} \punktebox{6}

{\small\textit{Erinnerung:} \es{conocer + a} bei Personen.}

\vspace{4pt}

\begin{teil}
    \item \es{¿Conoces} \luecke{1cm} \es{mi madre?}
    \item \es{No conozco} \luecke{1cm} \es{esta ciudad.}
    \item \es{Queremos conocer} \luecke{1cm} \es{tus amigos.}
    \item \es{¿Conocéis} \luecke{1cm} \es{el nuevo piso de Clara?}
    \item \es{Todavía no conozco} \luecke{1cm} \es{Nicolás.}
    \item \es{Mi hermano conoce} \luecke{1cm} \es{todas las canciones de este grupo.}
\end{teil}

\end{multicols}

% ═══════════════════════════════════════════════════════════════
% AUFGABE 4: Volle Breite (visuell abgesetzt)
% ═══════════════════════════════════════════════════════════════

\vspace{4pt}
\noindent\textcolor{hellgrau}{\rule{\textwidth}{1.5pt}}
\vspace{2pt}

\aufgabe{4}{Eigene Sätze.} \punktebox{7}

\begin{multicols}{2}

\textbf{a)} Beantworte mit \es{parecer}: \hfill{\small(3\,P.)}

\vspace{6pt}

\es{¿Qué te parece tu clase de español?}

\vspace{3pt}
\luecke{6cm}

\vspace{8pt}

\es{¿Qué te parecen los deberes?}

\vspace{3pt}
\luecke{6cm}

\columnbreak

\textbf{b)} Zwei Sätze mit \es{conocer}: \hfill{\small(4\,P.)}

{\small Achte auf \es{a} bei Personen!}

\vspace{6pt}
Person: \luecke{5.5cm}

\vspace{6pt}
Sache: \luecke{5.5cm}

\end{multicols}

% ═══════════════════════════════════════════════════════════════
% FOOTER
% ═══════════════════════════════════════════════════════════════

\vfill

\noindent\rule{\textwidth}{0.5pt}

\vspace{4pt}

\hfill\sanstext{\textbf{Punkte:}} \luecke{1.2cm} / \maxpunkte

\end{document}
